\section{The design decision}

\textbf{Build a proof coordinator around \grind, not a replacement for its
successful core.} The proposed tactic, \forge, uses a protected stock-search
lane and a richer search lane. The latter combines backward demands, checked
forward consequences, bounded structural proof search, and specialized
certificate producers. Crucially, those producers can return \emph{useful
lemmas}, not only a verdict on the entire goal. An inequality certificate can
supply an induction step; an induction lemma can expose arithmetic facts;
a synthesized witness can turn a quantified problem into a local \grind\ call.

``More powerful'' has three distinct meanings. \emph{Coverage} concerns which
proof patterns the tactic can discover. \emph{Efficiency} concerns resources
used on common goals. \emph{Reliability} concerns successful replay, predictable
failure, and the trusted assumptions of the resulting theorem. Forge targets
all three, but this report provides experimental evidence only for isolated
Python mechanisms. It does not claim measured Lean-level coverage gains or a
universal speedup.

\subsection{What should become easier}

The intended gains are concentrated where an otherwise short proof needs one
missing idea: strengthening an induction hypothesis, manufacturing a useful
term absent from the ground context, choosing a product of known nonnegative
expressions, finding an existential witness, or extracting a finite Boolean
subproblem. A large library can contain all the necessary lemmas and still
leave these decisions unresolved.

For example, from $x\geq 1$, $y\geq 2$, and $z\geq 3$, a useful intermediate
fact is
\[
 (x-1)(y-2)(z-3)\geq0.
\]
This is not difficult mathematics. The automation problem is selecting that
product, recording its dependencies, and making it available when a polynomial
normalizer or another solver needs it. Likewise, a proof about an accumulator
function may need a universally quantified accumulator invariant before any
amount of ground simplification helps.

The proposed architecture makes those choices explicit, bounded, inspectable,
and replayable. It does not promise automatic discovery of arbitrary creative
mathematics, complete higher-order proof search under finite resources, or
complete nonlinear arithmetic from a fixed-degree certificate basis.

\subsection{Implemented, specified, and still to be built}

\begin{center}
\begin{tabularx}{\textwidth}{@{}p{0.28\textwidth}YY@{}}
\toprule
Component & Delivered evidence & Production work still needed\\
\midrule
Demand-directed inference & Python tabling, dependency closure, proof DAG checker, differential tests & Typed Lean terms, e-class matching, library index, scoped metavariables\\
Nonlinear cone & Exact sparse identities; LP proposals; rational reconstruction & Lean reification, checker theorem or proof replay integration, scheduling\\
Bernstein positivity & Exact box certificates and complete-cover checking & Lean basis lemmas, dependent proof replay, richer split policies\\
Recurrence invariants & Interpolation proposals; exact base and transition identities & Extract Lean recursive definitions and apply appropriate recursors\\
Induction and witnesses & Detailed algorithms and worked derivations & Implement search, generalization, typed candidate synthesis\\
SAT/SMT/superposition & Concrete adapters and trust contracts specified & Version-compatible integrations, proof-rule coverage, testing\\
Whole \forge\ tactic & Architecture, interface, acceptance protocol & Not implemented; no Lean head-to-head results\\
\bottomrule
\end{tabularx}
\end{center}

\section{Begin with the actual baseline}

\subsection{Do not underestimate current \grind}

The current reference describes a cooperating collection of congruence closure,
constraint propagation, E-matching, guided case analysis, and theory solvers.
Its examples include nontrivial polynomial identities and finite-field
reasoning, not merely linear inequalities. It builds ordinary Lean proof
terms. The documentation also identifies large combinatorial Boolean search
as a poor fit and recommends a SAT-backed route for such problems
\cite{lean-grind}.

Consequently, ``add nonlinear arithmetic'' is too imprecise as a contribution.
Forge's algebraic extension is \emph{search for nonnegativity certificates and
useful inequality cuts}, rather than rediscovering polynomial equality
normalization. Similarly, ``add an e-graph,'' ``share facts,'' and ``support
plugins'' describe capabilities already present, not substantive extensions.

The inspected implementation and API also contain theorem guards, delayed
instances, support for extensionality and function congruence,
model-based theory-combination hooks, and a solver-extension interface.
The \code{Action} control layer uses continuations to accommodate branching,
nonchronological backtracking, and replay-script extraction
\cite{grind-types,grind-action}. Forge should exploit those mechanisms instead
of replacing them with an incompatible second notion of proof state.

\subsection{Related systems to reuse rather than relabel}

Aesop supplies a useful model for normalization rules and safe/unsafe
AND--OR search \cite{aesop}. Mathlib's \code{linarith} already separates an
untrusted search oracle from certificate reconstruction; \code{nlinarith}
adds nonlinear preprocessing \cite{linarith}. Thus, an LP oracle followed by an
exact identity check is not itself novel. The extra proposal is a goal-sensitive
product-and-square dictionary, cooperation with structural search, and the
protocol for feeding certified consequences back into a persistent context.

Duper is a concrete candidate for bounded superposition-based reasoning,
including higher-order capabilities \cite{duper}. Lean-SMT demonstrates a
route through preprocessing and proof-producing external SMT search, with
Lean-side reconstruction \cite{lean-smt}. These systems are optional lanes,
not reasons to erase type information or trust an external solver's answer.
A dependency must pass a compatibility and proof-reconstruction test before it
is enabled; the mere existence of a repository does not establish compatibility
with the target toolchain.

Recent work on interventions inside \grind\ is particularly relevant. It
reports that always-on changes can rescue some problems while breaking other
previously successful searches, and advocates interventions after baseline
failure \cite{interventions}. Forge adopts that deployment discipline. It does
not claim that a fallback cascade or bounded lookahead is a new invention.

\subsection{Version and provenance}

The inspected Lean release is \textbf{4.34.0}. The supplied replay project pins
mathlib to commit
\begin{center}
\code{1cf325a0cf67aca2b04d76b5380ff6a9e410aefa}.
\end{center}
The upstream toolchain at that revision specifies Lean 4.34.0. These facts were
read on 14 September 2026 in the America/Los\_Angeles time zone
\cite{lean-release,mathlib-revision}. The corresponding commit timestamp can
fall on 15 September in UTC. Pinning a compatible-looking pair is not a
successful build: the generated project was not compiled here.

\section{Architecture: proof obligations rather than opaque tactic calls}

\subsection{The persistent state}

At a logical level a search node is a sequent $\Gamma\vdash G$, together with a
resource account and a persistent snapshot. Its auxiliary data are:
\begin{itemize}
\item a typed term index and the local \grind\ state, with equality explanations;
\item certified facts, pending guarded facts, and backward demands;
\item an AND--OR graph of proof obligations and alternatives, with dependency
and scope information;
\item a certificate store, learned failure summaries, and a replay-cost estimate.
\end{itemize}

The distinction between \emph{fact} and \emph{demand} is fundamental. A fact
contains an expression $p$ and a proof of $p$ in a specified context. A demand
merely asks for such a proof. A proposed lemma is not inserted into the fact
store until it has been proved. Otherwise, two unsuccessful searches could
``prove'' each other through circular assumptions.

Represent a guarded proposal as
\[
 (\Gamma;\ g_1,\ldots,g_k;\ p;\ c),
 \qquad c:\Gamma\vdash g_1\to\cdots\to g_k\to p.
\]
The watchers for the $g_i$ may awaken when equalities or arithmetic bounds
change. Once each guard has a proof in a permitted scope, ordinary application
of $c$ produces a fact. Before that point, $p$ is only a conditional proposal.
This is also how a division normalization can wait for a nonzero denominator
instead of silently assuming it.

\subsection{Scope, identity, and transport}

Every stored theorem application records its local-context identifier,
free-variable dependencies, instantiated universes, typeclass instances,
metavariable assignments, and proof parents. A fact valid under a branch
assumption cannot leak into its sibling. To move a fact outside its original
scope, first abstract the branch assumptions and then prove the resulting
implication in the receiving context.

For example, a branch may prove $x^{-1}x=1$ using $x\ne0$. After backtracking,
the reusable result is $x\ne0\to x^{-1}x=1$, not the unguarded equality.
Caching only the printed expression loses this distinction. Environment
revision, transparency policy, and dependency closure must be part of a cache
key or explicitly revalidated on lookup.

Lean expressions also cannot be treated as an untyped term algebra.
Homogeneous equality requires compatible types; dependent applications may
need heterogeneous equality and transport. Existing \grind\ services should
supply equality explanations. Forge should never identify terms just because
their pretty-printed strings agree, or because an external first-order encoding
mapped them to the same symbol.

\subsection{The cooperative action protocol}

A plugin returns one of four logical outcomes: no applicable action, checked
progress, a closed proof, or an unresolved result with a reason. A progress
result carries a certified lemma or a well-formed refinement of the goal into
subgoals. Failure, numerical trouble, unsupported syntax, and timeout do not
mean falsehood.

\begin{Code}
cooperate(node, budget):
  saturate cheap certified consequences
  if the goal has a checked proof: return that proof
  enqueue demands extracted from the goal and pending guards
  repeat while resources remain:
    choose an action with a bounded local budget
    run it against an isolated snapshot
    if it returns a certificate:
      reconstruct and check its exact proposition
      merge only scoped, certified consequences
    if it returns subgoals:
      register an AND-node and its proof-combining continuation
    if it closes an OR-alternative:
      retain its proof and cancel unnecessary siblings
    age postponed actions; periodically service the fair lane
  return unknown with residual obligations and resource statistics
\end{Code}

This is stronger than repeatedly restarting \code{first | grind | nlinarith |
aesop}: the plugins exchange selected lemmas, guard proofs, and demand
information. Nevertheless, the first usable integration can begin with
independent leaf calls while the shared-state adapter is tested. That
engineering milestone should not be advertised as the complete architecture.

\subsection{A precise fallback guarantee}

Let $B(G,H)$ denote the result of stock \grind\ on $G$ under a fixed budget
$H$ and a fixed environment. Run that exact computation first. If it succeeds,
return its proof. Only if it fails, restore the initial snapshot and spend an
additional budget $K$ on Forge's extensions.

\begin{proposition}[Additive-budget preservation]
Assuming identical baseline inputs, resource accounting, and execution
conditions, the cascade succeeds on every goal for which $B(G,H)$ succeeds.
\end{proposition}
\begin{proof}
The successful baseline execution is performed unchanged, and its result is
returned before any extension is run.
\end{proof}

This elementary guarantee is about \emph{additional resources}. It is not a
proof that allocating part of the same total budget $H$ to new machinery
preserves all baseline successes. Nor does a simultaneously running baseline
thread receive an identical wall-clock or memory budget if the other threads
contend with it. Report additive-budget and fixed-total-budget experiments
separately. Fairness prevents starvation in an unbounded search; it does not
make a bad finite-budget allocation free.

\section{Demand-directed instantiation and intermediate lemmas}

\subsection{Backward demands with forward proof completion}

Suppose a theorem has an instantiated conclusion matching the wanted fact:
\[
 A_1(\vec t)\to\cdots\to A_k(\vec t)\to C(\vec t).
\]
Rather than forward-generating every instance of every theorem, create demands
for the $A_i(\vec t)$ when $C(\vec t)$ is useful. A conclusion index retrieves
candidate declarations by typed head symbols, argument structure, and relevant
relation classes. Matching is modulo certified equality when useful, but
elaboration decides whether the proposed instantiation is legal.

A candidate whose conclusion does not determine every parameter needs a second
phase: enumerate a bounded typed term grammar, solve a small witness problem,
or leave it to the fair forward lane. It must not arbitrarily instantiate a
universally quantified theorem variable with an ill-scoped metavariable.

The Python prototype deliberately avoids that difficulty. Its rules are
single-sort Horn clauses whose body variables all occur in the head. Matching
a \emph{ground} demanded head therefore determines the substitution for every
premise. There is no higher-order unification, dependent typing, equality
reasoning, or library retrieval in this prototype.

\subsection{Tabling prevents cyclic pseudo-proofs}

Build a finite, budgeted backward closure of ground demands and candidate rule
instances. Then propagate actual proofs forward over that dependency graph.
Each rule instance has a counter of unresolved premises. Input facts initialize
the agenda; proving a premise decreases the counters of its dependents. A rule
fires only when every premise is proved. Store the proof node at that moment.

This is not a depth-first convention that treats ``currently being visited''
as success. Rules $P\to Q$ and $Q\to P$, without a base fact, prove neither.
The certificate is a topologically ordered DAG: every premise index is strictly
smaller than its conclusion's index. Its checker reconstructs the exact
substitution and conclusion for each theorem application.

\begin{proposition}[Relative completeness of the finite Horn core]
For a fixed finite set of ground candidate rule instances, if the dependency
agenda is exhausted without a resource cutoff, every fact in their least
forward closure is eventually proved. In particular, a backward-relevant
closure containing a ground derivation of the target will discover a proof.
\end{proposition}
\begin{proof}
Induct on the height of a finite derivation. Height-zero input facts enter the
agenda initially. Once all smaller-height premises of an instance have entered,
its unresolved counter reaches zero, so its conclusion is enqueued. Tabling
ensures repeated discoveries do not obstruct this propagation.
\end{proof}

The statement concerns the fixed finite closure, not unrestricted first-order
logic. A depth cap can exclude a necessary derivation. With function symbols,
the full relevant closure may be infinite. Such outcomes are reported as
\code{unknown}, not as refutations.

\subsection{A family that isolates the advantage}

Take a base fact $P(a)$ and the two rules
\[
 P(x)\to P(g(x)),\qquad P(x)\to P(f(x)).
\]
Ask for $P(f^d(a))$. In the supplied breadth-first forward implementation, the
$g$-rule precedes the $f$-rule. The goal is therefore the last word generated at
depth $d$, after
\[
 1+2+\cdots+2^d=2^{d+1}-1
\]
facts. Backward demand generation produces exactly the $d+1$ facts on the
$f$-chain. Both returned proofs contain $d+1$ nodes.

At depth 14 the measured counts are 32,767 versus 15. This is an intentionally
adversarial \emph{Python Horn experiment}, not a reconstruction of \grind.
Current \grind\ already sees a negated goal and performs sophisticated
instantiation; this example cannot establish that it would suffer the same
expansion. What the experiment establishes is the implemented dependency
mechanism and its exact work reduction on this family.

\subsection{Library retrieval and induction of useful demands}

A production index should first use deterministic structure: conclusion head,
constants, type constructor, theorem arity, and normalized argument positions.
An inexpensive text ranker can supplement these features. Neural ranking is
optional; it is never proof evidence. After retrieving a declaration, instantiate
it through Lean's elaborator, generate its side goals, and attach its actual
proof term. Never trust a textual theorem statement supplied by a model.

Useful demands need not be the final goal. Sources include a missing denominator
condition, the guard of a conditional rewrite, the precondition of an induction
step, and a sign condition requested by the nonlinear solver. Give these
consumers a small number of ranked requests rather than internalizing every
retrieved theorem globally. Cache failed instantiations by the exact context
and budget, but invalidate failures when relevant facts change.

Demand search, like any relevance filter, can miss a surprising intermediate
lemma. Reserve a fair lane for ordinary instantiation and gradually widen the
retrieval and term-generation bounds. The practical objective is to find a
short proof cheaply, not to turn a relevance heuristic into an unjustified
completeness claim.
